\documentclass[aspectratio=169,11pt]{beamer}

\usepackage{ubc-moloch-beamer}

% Show note text in the exported PDF. Remove this line if you only want slides.
\setbeameroption{show notes on second screen}
\setbeamertemplate{note page}[plain]

\author{Style API Matrix}

\newcommand{\DocSetTitleMeta}[1]{%
  \title{ubc-moloch-beamer.sty}%
  \subtitle{#1}%
  \institute{Theme behavior matrix}%
  \date{\today}%
}

\newcommand{\DocExplainAfter}[2]{%
  \begin{frame}
    \frametitle{#1}
    The previous slide is the real style case for this configuration.
    \note{#2}
  \end{frame}
}

\begin{document}

\UBCFrameSetup{crest=off}
\begin{frame}
  \frametitle{How to read this file}
  This `style-matrix.tex` is a usage document for `ubc-moloch-beamer.sty`.
  It is based on the moloch beamer theme \cite{jolars-moloch}.
  \begin{itemize}
    \item The deck uses only \texttt{\string\UBCSetup}, \texttt{\string\UBCFrameSetup}, and \texttt{\string\UBCTitlePage}.
    \item Title and section cases are documented by the slide that follows them.
    \item Regular-frame cases carry their explanation in \texttt{\string\note} placed directly in the frame body.
  \end{itemize}
  \note{
    This deck intentionally ignores thesis content and only demonstrates style behavior.
    Case keys:
    - T = title.
    - S = section.
    - N = normal.
    - U = standout.
    Title pages now expose only the title variant.
    Standout frames are always dark and crest-free by package design.
  }
\end{frame}

% ============================================================================
% 1) TITLE FRAME CASES (2)
% ============================================================================

\DocSetTitleMeta{Case T-L: Title page, light}
\UBCTitlePage[variant=light]
\DocExplainAfter{Case T-L (Title)}{
  Case T-L.
  Input:
  - UBCTitlePage[variant=light]
  Expected:
  - light title styling
  - corner crest hidden
  - titlegraphic defaults to \string\ubccrest when unset
}

\DocSetTitleMeta{Case T-D: Title page, dark}
\UBCTitlePage[variant=dark]
\DocExplainAfter{Case T-D (Title)}{
  Case T-D.
  Input:
  - UBCTitlePage[variant=dark]
  Expected:
  - dark title styling
  - corner crest hidden
  - titlegraphic defaults to \string\ubccrest when unset
}

% ============================================================================
% 2) SECTION FRAME CASES (4)
% ============================================================================

\UBCSetup{section variant=light, section crest=off}
\section{Case S-L-OFF}
\DocExplainAfter{Case S-L-OFF (Section)}{
  Case S-L-OFF.
  Input before \string\section{...}:
  - UBCSetup{section variant=light, section crest=off}
  Expected section page:
  - light variant
  - corner crest hidden
}

\UBCSetup{section variant=light, section crest=on}
\section{Case S-L-ON}
\DocExplainAfter{Case S-L-ON (Section)}{
  Case S-L-ON.
  Input before \string\section{...}:
  - UBCSetup{section variant=light, section crest=on}
  Expected section page:
  - light variant
  - corner crest visible
}

\UBCSetup{section variant=dark, section crest=off}
\section{Case S-D-OFF}
\DocExplainAfter{Case S-D-OFF (Section)}{
  Case S-D-OFF.
  Input before \string\section{...}:
  - UBCSetup{section variant=dark, section crest=off}
  Expected section page:
  - dark variant
  - corner crest hidden
}

\UBCSetup{section variant=dark, section crest=on}
\section{Case S-D-ON}
\DocExplainAfter{Case S-D-ON (Section)}{
  Case S-D-ON.
  Input before \string\section{...}:
  - UBCSetup{section variant=dark, section crest=on}
  Expected section page:
  - dark variant
  - corner crest visible
}

% ============================================================================
% 3) NORMAL FRAME CASES (4)
% ============================================================================

\UBCFrameSetup{variant=light, crest=off}
\begin{frame}
  \frametitle{Case N-L-OFF}
  Normal frame, light variant, crest request OFF.
  \note{
    Case N-L-OFF.
    Input:
    - UBCFrameSetup{variant=light, crest=off}
    Expected:
    - light body background
    - corner crest hidden
  }
\end{frame}

\UBCFrameSetup{variant=light, crest=on}
\begin{frame}
  \frametitle{Case N-L-ON}
  Normal frame, light variant, crest request ON.
  \note{
    Case N-L-ON.
    Input:
    - UBCFrameSetup{variant=light, crest=on}
    Expected:
    - light body background
    - corner crest visible
  }
\end{frame}

\UBCFrameSetup{variant=dark, crest=off}
\begin{frame}
  \frametitle{Case N-D-OFF}
  Normal frame, dark variant, crest request OFF.
  \note{
    Case N-D-OFF.
    Input:
    - UBCFrameSetup{variant=dark, crest=off}
    Expected:
    - dark body background
    - corner crest hidden
  }
\end{frame}

\UBCFrameSetup{variant=dark, crest=on}
\begin{frame}
  \frametitle{Case N-D-ON}
  Normal frame, dark variant, crest request ON.
  \note{
    Case N-D-ON.
    Input:
    - UBCFrameSetup{variant=dark, crest=on}
    Expected:
    - dark body background
    - corner crest visible
  }
\end{frame}

% ============================================================================
% 4) STANDOUT FRAME CASES (2)
% ============================================================================

\UBCSetup{normal crest=off}
\begin{frame}[standout]
  Case U-OFF
  \note{
    Case U-OFF.
    Input:
    - UBCSetup{normal crest=off}
    Expected:
    - dark standout style
    - corner crest hidden
  }
\end{frame}

\UBCSetup{normal crest=on}
\begin{frame}[standout]
  Case U-ON
  \note{
    Case U-ON.
    Input:
    - UBCSetup{normal crest=on}
    Expected:
    - same visual result as U-OFF
    - standout frames keep the corner crest hidden by design
  }
\end{frame}

\begin{frame}[allowframebreaks]
  \frametitle{References}
  \bibliographystyle{plain}
  \bibliography{references}
\end{frame}

\end{document}
